% 致谢
\begin{acknowledgement}
  时光荏苒，研究生阶段的学习与科研工作即将画上句号。在论文完成之际，谨向所有关心、帮助和支持过我的老师、同学、亲友致以最诚挚的感谢。

  首先，衷心感谢我的导师袁建军老师。袁老师为本课题提供了宝贵的实验资源与研究平台，并在课题方向把握、研究思路凝练以及学术规范养成等方面给予了我持续而细致的指导，使我受益匪浅。

  同时，感谢实验室各位指导老师的悉心帮助，特别感谢杜亮老师、鲍晟老师在科研训练和日常学习中给予我的指导与支持。尤其要特别感谢胡正涛老师，在论文撰写、实验分析和细节打磨过程中给予了我大量耐心、专业且富有启发性的指导，使论文得以不断完善。

  此外，衷心感谢南京大学张振宇老师、福建农林大学陈昊老师对本论文提出的宝贵指导与中肯建议。两位老师从不同视角提出了许多建设性意见，为论文内容的完善和研究深度的提升提供了重要帮助。

  感谢实验室的师兄师姐、师弟师妹们在科研与生活中的支持与帮助。与你们共同学习、讨论和合作的经历，让我在面对困难时始终能够获得力量，也让我度过了充实而难忘的研究生时光。

  最后，感谢我的家人和朋友长期以来的理解、包容与鼓励。正是你们无条件的支持，让我能够心无旁骛地投入学习和研究，并坚持走到今天。

  谨以此文，向所有给予我帮助的人致以诚挚谢意！

  \vspace{50bp}

  \hfill 作者署名：邱模聪

  \hfill 完成地点：上海大学

  \hfill 2026 年 4 月 1 日

\end{acknowledgement}
